% TEMPLATE-MILC-v3.tex - plantilla maestra UNICA de las guias MILC v3.
%
% La usan las cuatro familias de guias, cada una con su driver:
%   · Grados 6-11          -> scripts/build-guias-g11.py
%   · Bebras               -> scripts/build-guias-bebras.py
%   · Semillero            -> scripts/build-guias-semillero.py
%   · Territorio interior  -> scripts/build-guias-territorio-interior.py
%
% Antes habia tres copias casi identicas (~400 lineas iguales, 6 distintas):
% cada mejora habia que aplicarla tres veces, y en la practica no se hacia
% (el motor de recursos vivia solo en dos de ellas). Lo que varia entre
% familias esta parametrizado en 6 placeholders que cada driver llena
% (se nombran aqui SIN su sintaxis real: el builder reemplaza por texto plano,
% tambien dentro de los comentarios, y un valor multilinea rompe el preambulo):
%
%   HEADER_IZQ        encabezado izquierdo de pagina
%   HEADER_DER        encabezado derecho de pagina
%   PORTADA_ETIQUETA  rotulo sobre el numero grande (GRADO/MOMENTO/MODULO)
%   PORTADA_SERIE     linea de serie de la portada
%   PORTADA_ID        identificador de la guia en la portada
%   PORTADA_CIRCULO   el circulito de grado (°), o vacio si no aplica
%   PDF_TITULO        titulo en texto plano (sin \textbf ni comillas TeX) para
%                     los metadatos del PDF (/Title); lo produce lib_xelatex.texto_pdf
%
% Composicion (auditoria 2026-09-03): las bandas de estacion piden espacio
% (\Needspace*) y prohiben el salto justo despues (\nopagebreak), asi no quedan
% huerfanas al pie; las cajas largas (softbox, drawbox, infoband, puentebox) son
% `breakable`, asi una caja de media pagina ya no arrastra todo a la siguiente
% dejando la anterior al 40 %. Todo texto sobre mostaza o verde va en milcNegro
% (blanco sobre #E5B400 da 1,93:1 y sobre #58A923 2,95:1; ninguno pasa WCAG AA).
% El resto de claves las documenta content/guias/_SCHEMA.md. El script valida
% que no queden placeholders sin reemplazar antes de escribir el archivo final.

% Etiquetado PDF (latex-lab): /Lang, estructura y texto alternativo de las
% figuras, para lectores de pantalla. Debe ir ANTES de \documentclass.
\DocumentMetadata{lang=es-CO,tagging=on}
\documentclass[11pt,letterpaper]{article}
% headheight=24pt: la cabecera derecha lleva dos lineas (identidad de la guia
% y estacion actual). headsep se acorta en la misma medida para que el bloque
% de texto quede exactamente donde estaba con headheight=12pt.
\usepackage[margin=2.0cm,headheight=24pt,headsep=13pt]{geometry}
\usepackage[table]{xcolor}
\usepackage{fontspec}
\usepackage[spanish]{babel}
\usepackage{tikz}
% Librerías TikZ para los diagramas de las guías (bloque `recursos:`):
%  · babel  → babel en español vuelve activos caracteres como «>», y TikZ los usa
%             en sus opciones (>=stealth, ->). Sin esta librería, cualquier
%             diagrama con flechas revienta con «\language@active@arg> has an
%             extra }». Debe ir ANTES de que se dibuje nada.
%  · positioning        → sintaxis «below left=2pt and 3pt».
%  · decorations.pathreplacing → llaves ({brace}) para señalar rangos.
\usetikzlibrary{babel,positioning,decorations.pathreplacing}
\usepackage{graphicx}
% Circuitos: el trabajo de electrónica se hace hoy con micro:bit y sus sensores
% integrados (MakeCode, sin cableado), así que no se necesitan esquemáticos.
% Si algún día entran montajes con protoboard, basta descomentar esta línea:
% los diagramas .tex/.tikz declarados en `recursos:` ya se incrustan solos.
% \usepackage{circuitikz}
\usepackage[most]{tcolorbox}
\usepackage{tabularx}
\usepackage{array}
\usepackage{fancyhdr}
\usepackage{titlesec}
\usepackage[document]{ragged2e}  % bandera a la derecha en todo el cuerpo
\usepackage{enumitem}
\usepackage{microtype}
\usepackage{etoolbox}
\usepackage{needspace}
% hyperref al final del preambulo: metadatos del PDF (titulo, autor, idioma,
% familia), marcadores por estacion (\pdfbookmark) y enlaces sin recuadro.
\usepackage[hidelinks,
  pdftitle={Iterar y refinar texto con IA — del borrador al capítulo terminado},
  pdfauthor={PhD. Álvaro Cárdenas Orozco},
  pdfsubject={Tecnología e Informática · Grado 10 · Guía 4 · MILC v3},
  pdflang={es-CO},
  bookmarksopen=true,bookmarksnumbered=false]{hyperref}

% Helvetica Neue no trae la flecha U+2192, asi que cada «→» escrito literalmente
% en el YAML se compilaba a NADA, en silencio (462 flechas en 61 guias: rutas del
% tipo «paleta → tipografia → lockup» salian rotas). Mapeamos el caracter a su
% version matematica, que si tiene glifo. Asi el autor puede seguir escribiendo
% «→» comodamente en el YAML.
\usepackage{newunicodechar}
\newunicodechar{→}{\ensuremath{\rightarrow}}
% Mismo problema con los simbolos de marca y vinetas: el «✓» de «una palomita ✓»
% salia como cuadrito vacio en el PDF de todas las guias que lo usaban.
\usepackage{pifont}
\newunicodechar{✓}{\ding{51}}
\newunicodechar{✗}{\ding{55}}
\newunicodechar{✔}{\ding{52}}
\newunicodechar{•}{\textbullet}
\newunicodechar{≥}{\ensuremath{\geq}}
\newunicodechar{≤}{\ensuremath{\leq}}

\setmainfont{Helvetica Neue}
\setsansfont{Helvetica Neue}
\setlength{\parindent}{0pt}
\setlength{\parskip}{6pt}
% Sin particion de palabras: en 6.º-7.º una palabra cortada por guion al final
% de linea («Comuni-cacion») frena la lectura mucho mas que una linea corta.
\hyphenpenalty=10000
\exhyphenpenalty=10000
\pretolerance=2000
\tolerance=3000
\setlength{\emergencystretch}{3em}
\linespread{1.10}
\renewcommand{\arraystretch}{1.25}
\setlist{leftmargin=5mm,itemsep=1mm,topsep=1mm}
\newcolumntype{Y}{>{\RaggedRight\arraybackslash}X}

\definecolor{milcMagenta}{HTML}{E6007E}
\definecolor{milcVino}{HTML}{5A0038}
\definecolor{milcTurquesa}{HTML}{0093A5}
\definecolor{milcVerde}{HTML}{58A923}
\definecolor{milcMostaza}{HTML}{E5B400}
\definecolor{milcOcre}{HTML}{B86600}
\definecolor{milcNegro}{HTML}{241816}
\definecolor{milcGris}{HTML}{F7F7F4}
\definecolor{milcLinea}{HTML}{E8E2DD}
\definecolor{milcGradePrimary}{HTML}{FF2D87}
\definecolor{milcGradeDark}{HTML}{9F1B56}
\definecolor{milcGradeSoft}{HTML}{FFE0EE}
\definecolor{milcGradeAccent}{HTML}{CFE7FF}
\definecolor{milcGradeText}{HTML}{FFFFFF}
\color{milcNegro}

\pagestyle{fancy}
\fancyhf{}
% Cabecera de dos lineas: identidad de la guia arriba y, a la derecha y en
% gris, la estacion en curso (\markright desde \milcestacion). Antes de la
% primera estacion la segunda linea va vacia; el \strut mantiene la altura
% para que la linea de arriba no baile entre paginas.
\lhead{\small Tecnología e Informática · Grado 10\\[-1pt]{\footnotesize\strut}}
\rhead{\small Guía 4 · MILC v3\\[-1pt]{\footnotesize\strut\color{milcNegro!65}\rightmark}}
\cfoot{\small\thepage}
\renewcommand{\headrulewidth}{0pt}

% Sobre mostaza (#E5B400) y verde (#58A923) el texto va en milcNegro; sobre el
% resto de la paleta, en blanco. Se usa dentro de fonttitle/fontupper, donde un
% \color posterior manda sobre coltitle.
\newcommand{\milctintasobre}[1]{%
  \ifboolexpr{test{\ifstrequal{#1}{milcMostaza}} or test{\ifstrequal{#1}{milcVerde}}}%
    {\color{milcNegro}}{\color{white}}}

\newtcolorbox{titlebox}[1]{
  enhanced,colback=#1,colframe=#1,arc=7mm,boxrule=0pt,
  left=7mm,right=7mm,top=3mm,bottom=3mm,
  before skip=8pt,after skip=8pt,
  fontupper=\sffamily\bfseries\Large\color{white}
}

\newtcolorbox{softbox}[3]{
  enhanced,breakable,pad at break=3mm,
  colback=#2,colframe=#1,borderline west={4pt}{0pt}{#1},
  arc=2mm,boxrule=.4pt,left=4mm,right=4mm,top=3mm,bottom=3mm,
  before skip=6pt,after skip=8pt,title={#3},coltitle=white,
  colbacktitle=#1,fonttitle=\sffamily\bfseries\milctintasobre{#1},fontupper=\RaggedRight,
  attach boxed title to top left={xshift=3mm,yshift=-2mm},
  boxed title style={arc=2mm, boxrule=0pt}
}

\newtcolorbox{drawbox}[2]{
  enhanced,breakable,pad at break=4mm,
  colback=white,colframe=#1,arc=4mm,boxrule=1pt,
  left=5mm,right=5mm,top=4mm,bottom=4mm,before skip=8pt,after skip=8pt,
  title={#2},coltitle=white,colbacktitle=#1,fonttitle=\sffamily\bfseries\milctintasobre{#1},
  fontupper=\RaggedRight,
  attach boxed title to top left={xshift=4mm,yshift=-2mm},
  boxed title style={arc=3mm, boxrule=0pt}
}

\newtcolorbox{infoband}[1]{
  enhanced,breakable,pad at break=2mm,colback=#1!8,colframe=#1!50,arc=2mm,boxrule=.5pt,
  left=4mm,right=4mm,top=2mm,bottom=2mm,before skip=4pt,after skip=4pt,
  fontupper=\small\RaggedRight
}

% Puente narrativo entre secciones (contrato editorial Conéctate).
% Da continuidad: "Acabas de X. Ahora vas a Y porque Z."
% Visualmente: banda izquierda en color del grado, texto en itálica pequeña.
\newtcolorbox{puentebox}{
  enhanced,breakable,pad at break=2mm,colback=milcGradeSoft,colframe=milcGradeDark,
  borderline west={3pt}{0pt}{milcGradeDark},
  arc=0mm,boxrule=0pt,left=5mm,right=4mm,top=2.5mm,bottom=2.5mm,
  before skip=4pt,after skip=8pt,
  fontupper=\itshape\small\color{milcNegro}\RaggedRight
}
% La flecha va en modo matematico a proposito: Helvetica Neue NO tiene el glifo
% U+2192, asi que \textrightarrow se compilaba a NADA («Missing character: There
% is no → in font Helvetica Neue») y el puente salia sin flecha. En modo
% matematico la flecha viene de la fuente matematica y si se dibuja.
\newcommand{\puente}[1]{%
  \begin{puentebox}%
    \textcolor{milcGradeDark}{$\rightarrow$}\hspace{2mm}#1%
  \end{puentebox}%
}

\newcommand{\linead}[1]{\noindent\color{milcLinea}\rule{#1}{0.5pt}\color{milcNegro}}
\newcommand{\respuesta}[1]{\par\vspace{2mm}\foreach \n in {1,...,#1}{\noindent\color{milcLinea}\rule{\linewidth}{0.5pt}\par\vspace{5mm}}\color{milcNegro}}
\newcommand{\checkbox}{\textcolor{milcGradeDark}{\fbox{\phantom{x}}}\hspace{5pt}}

% ─── Listas aireadas para las guías socioemocionales ────────────────────
% Componer las verificaciones como «\checkbox texto\\» deja el texto que salta
% de línea debajo del cuadrito y apelmaza el bloque. Estos entornos alinean la
% continuación y airean los ítems. Los usa Territorio Interior; cualquier
% familia puede hacerlo.
\newlist{checklista}{itemize}{1}
\setlist[checklista]{
  label=\textcolor{milcGradeDark}{\fbox{\phantom{x}}},
  leftmargin=9mm, labelsep=3.2mm, itemsep=2.4mm,
  topsep=2.5mm, parsep=0pt, partopsep=0pt, before=\RaggedRight
}
\newlist{pasoslista}{enumerate}{1}
\setlist[pasoslista]{
  label=\textcolor{milcGradeDark}{\sffamily\bfseries\arabic*.},
  leftmargin=8mm, labelsep=3mm, itemsep=2.8mm,
  topsep=1mm, parsep=0pt, partopsep=0pt, before=\RaggedRight
}

\newcommand{\phasecard}[4]{
\begin{tcolorbox}[enhanced,colback=#3,colframe=#2,arc=3mm,boxrule=.5pt,height=3.05cm,valign=center,halign=center,left=1.5mm,right=1.5mm,top=2mm,bottom=2mm]
{\sffamily\bfseries\fontsize{22}{24}\selectfont\textcolor{#2}{#1}}\par
{\sffamily\bfseries\small #4}
\end{tcolorbox}
}

% ── Piezas del rediseno 2026 (legibilidad 6.º-11.º) ───────────────────
% Banda de estacion: reemplaza el titlebox macizo. Titulo a la izquierda,
% dato practico (tiempo/modalidad) a la derecha, en una sola linea.
% Nunca huerfana: pide 9 lineas libres (o salta de pagina rellenando con
% \vfill) y prohibe el salto justo despues de la banda. Nueve y no seis:
% la caja partible que sigue necesita ~80pt (titulo adosado, relleno y dos
% lineas) para arrancar en la misma pagina; con menos, tcolorbox la manda
% entera a la siguiente y la banda se queda sola. El penalty 10000 va
% en `after`, ANTES del espacio posterior: un `after skip` (aunque sea 0pt)
% es un \vskip tras la caja, y ese glue es un punto de corte legal que un
% \nopagebreak escrito despues de \end{tcolorbox} ya no cierra. Cada banda
% deja marcador PDF y actualiza la cabecera derecha.
\newcounter{milcestacion}
\newcommand{\milcestacion}[2]{%
\Needspace*{9\baselineskip}%
\stepcounter{milcestacion}%
\pdfbookmark[1]{#2}{milcestacion\themilcestacion}%
\markright{#2}%
\begin{tcolorbox}[enhanced,colback=#1,colframe=#1,arc=2mm,boxrule=0pt,
  left=5mm,right=5mm,top=2.6mm,bottom=2.6mm,before skip=11pt,
  after={\par\nopagebreak\vskip7pt}]
{\sffamily\bfseries\fontsize{15}{18}\selectfont\milctintasobre{#1}#2}
\end{tcolorbox}}

% Tarjeta de paso numerada: sustituye las tablas «Pilar / Aplicacion», que
% eran formato de consulta docente, no de estudiante. El numero va en un
% circulo sobre el borde para que la secuencia se lea de un vistazo.
\newcommand{\milcpaso}[3]{%
\Needspace{4\baselineskip}%
\begin{tcolorbox}[enhanced,colback=white,colframe=milcLinea,arc=1.5mm,boxrule=.9pt,
  left=14mm,right=4mm,top=3mm,bottom=3mm,before skip=4pt,after skip=4pt,
  fontupper=\RaggedRight,
  overlay={\node[anchor=north west,circle,fill=#2,text=white,inner sep=0pt,
    minimum size=7.4mm,font=\sffamily\bfseries\small]
    at ([xshift=3.6mm,yshift=-3.2mm]frame.north west) {#1};}]
#3
\end{tcolorbox}}

% Casilla marcable de tamano util con lapiz (la anterior era \fbox{\phantom{x}},
% de alto variable segun el texto que la seguia).
\newcommand{\milcmarca}{\framebox[3.8mm]{\rule{0pt}{3.4mm}}}

% Fila marcable de lista suelta (checks de la estacion 1).
\newcommand{\milccheckrow}[1]{%
\par\vspace{1.8mm}\noindent
\makebox[6.4mm][l]{\raisebox{-.55mm}{\textcolor{milcNegro}{\milcmarca}}}%
\parbox[t]{\dimexpr\linewidth-6.4mm\relax}{\RaggedRight #1}\par}

% Celda marcable dentro de la rubrica (tabla de autoevaluacion). Misma
% sangria francesa, pero pensada para vivir dentro de una columna X.
\newcommand{\milcopcion}[1]{%
\makebox[5.2mm][l]{\raisebox{-.55mm}{\textcolor{milcNegro}{\milcmarca}}}%
\parbox[t]{\dimexpr\linewidth-5.2mm\relax}{\RaggedRight #1}}

% ── Recursos visuales de la guía ──────────────────────────────────────
% Declarados en el bloque `recursos:` del YAML y emitidos por el builder.
% \guiaFigura   → imagen rasterizada o PDF (foto, carta celeste, montaje).
% \guiaDiagrama → fuente TikZ/circuitikz (.tex) incrustada como vector:
%                 es la vía para circuitos y esquemáticos editables.
% [#1] = texto alternativo (alt) · #2 = ruta relativa al .tex · #3 = pie
% El alt llega al PDF etiquetado (/Alt) y lo lee el lector de pantalla.
\newcommand{\guiaFigura}[3][]{%
\begin{center}
\includegraphics[width=0.88\linewidth,height=0.38\textheight,keepaspectratio,alt={#1}]{#2}\\[1.5mm]
{\footnotesize\itshape\textcolor{milcNegro!75}{#3}}
\end{center}
}

\newcommand{\guiaDiagrama}[2]{%
\begin{center}
\input{#1}\\[1.5mm]
{\footnotesize\itshape\textcolor{milcNegro!75}{#2}}
\end{center}
}

\begin{document}
\thispagestyle{empty}

% ── Portada ───────────────────────────────────────────────────────────
% Antes: un rectangulo de color a pagina completa. Se veia bien en pantalla,
% pero en la fotocopiadora del colegio se comia el toner y salia gris sucio.
% Ahora la banda ocupa el tercio superior y el resto va en blanco.
\begin{tikzpicture}[remember picture,overlay]
  \path[fill=milcGradePrimary]
    (current page.north west) rectangle ([yshift=-10.6cm]current page.north east);

  \node[anchor=north west,text=milcGradeText] at ([xshift=2.0cm,yshift=-1.55cm]current page.north west)
    {\sffamily\bfseries\fontsize{12}{14}\selectfont GRADO};
  \node[anchor=north west,text=milcGradeText,opacity=.85] at ([xshift=2.0cm,yshift=-2.35cm]current page.north west)
    {\sffamily\bfseries\fontsize{8.5}{10}\selectfont SERIE GUÍAS · TECNOLOGÍA E INFORMÁTICA};
  \node[anchor=north east,text=milcGradeText,opacity=.85,align=right] at ([xshift=-2.0cm,yshift=-1.55cm]current page.north east)
    {\sffamily\bfseries\fontsize{9}{12}\selectfont I.E. Sor María Juliana\\Cartago, Valle del Cauca};

  \node[anchor=west,text=milcGradeText] (gradonum) at ([xshift=1.85cm,yshift=-7.1cm]current page.north west)
    {\sffamily\bfseries\fontsize{112}{112}\selectfont 10};
  % El circulito de grado (°) solo aplica a las guias de grado («11°»); en
  % Bebras y Semillero el numero grande es un momento/modulo, no un grado.
  % Se posiciona relativo al nodo `gradonum`, no en coordenadas absolutas.
  \draw[milcGradeText,line width=1.6pt]
    ([xshift=2.0mm,yshift=-4.5mm]gradonum.north east) circle (.15cm);

  \node[anchor=west,text=milcGradeText,text width=11.4cm,align=left] at ([xshift=1.5cm,yshift=0.5cm]gradonum.east)
    {
     {\sffamily\bfseries\fontsize{9}{11}\selectfont\MakeUppercase{Período 1 · Escritura abierta con IA}}\\[3mm]
     {\sffamily\bfseries\fontsize{21}{25}\selectfont\setlength{\baselineskip}{27pt}Iterar y refinar ---\\[4pt]del borrador al\\[4pt]capítulo terminado}\\[3.5mm]
     {\sffamily\bfseries\fontsize{15}{18}\selectfont Guía 4-10-TIC}
    };
\end{tikzpicture}

\vspace*{9.4cm}
\vfill

{\sffamily\bfseries\fontsize{8}{10}\selectfont\textcolor{milcGradeDark}{LO QUE TE LLEVAS HOY}}

\begin{tcolorbox}[enhanced,colback=white,colframe=milcNegro,arc=2mm,boxrule=1.6pt,
  left=5mm,right=5mm,top=4mm,bottom=4mm,before skip=3pt,after skip=12pt,
  fontupper=\RaggedRight\fontsize{11}{15.5}\selectfont]
Un capítulo en 3 versiones progresivas, con bitácora de cambios, comparación lado a lado y una reflexión honesta sobre cuánto del texto final es tuyo
\end{tcolorbox}

\begin{tcolorbox}[enhanced,colback=milcGris,colframe=milcGris,arc=2mm,boxrule=0pt,
  left=5mm,right=5mm,top=3.5mm,bottom=3.5mm,after skip=12pt]
{\sffamily\bfseries\fontsize{8}{10}\selectfont\textcolor{milcNegro!55}{ESTA GUÍA ES DE}}\\[3mm]
\begin{tabularx}{\linewidth}{X p{3.2cm} p{3.2cm}}
\linead{\linewidth} & \linead{3.0cm} & \linead{3.0cm}\\[-1mm]
{\fontsize{8.5}{10}\selectfont\textcolor{milcNegro!55}{Nombre completo}} &
{\fontsize{8.5}{10}\selectfont\textcolor{milcNegro!55}{Grupo}} &
{\fontsize{8.5}{10}\selectfont\textcolor{milcNegro!55}{Fecha}}\\
\end{tabularx}
\end{tcolorbox}

\vfill

\noindent\textcolor{milcLinea}{\rule{\linewidth}{1pt}}\\[2mm]
{\sffamily\bfseries\fontsize{9.5}{12}\selectfont Autor: Dr. Álvaro Cárdenas Orozco}\\[1mm]
{\sffamily\fontsize{8.5}{11}\selectfont\textcolor{milcNegro!55}{Metodología MILC · apertura ancestral · pensamiento computacional · triángulo Dussel–estoicismo–Floridi}}

\newpage

\section*{Apertura: Iterar y refinar texto con IA --- del borrador al capítulo terminado}

\begin{softbox}{milcGradeDark}{milcGradeSoft}{Saber ancestral}
En el norte del Valle hay un oficio hermano del bordado que trabaja quitando en vez de poniendo: el \textbf{calado}. Ansermanuevo se llama a sí misma cuna del calado y el bordado, y lo celebra cada agosto en sus fiestas (Soy Valle, s. f.). La técnica base es el deshilado. Se retiran hilos de la trama o de la urdimbre de la tela y se reagrupan con puntadas los que quedan, hasta que el vacío mismo dibuja. Hay oficios enteros, entonces, donde parte del trabajo consiste en deshacer. Y donde lo que se ve por delante depende de lo que se hizo por detrás. La \textbf{cara de exclusión} es la misma del bordado. Es un oficio precario, y el reconocimiento patrimonial en curso no equivale a un ingreso digno para las artesanas.\par\smallskip{\footnotesize\textcolor{milcNegro!70}{\textbf{Fuente.} Soy Valle. (s.\,f.). \emph{Ruta del bordado y el calado: Ansermanuevo, Cartago}. Consultado en 2026.}}
\end{softbox}

\begin{softbox}{milcTurquesa}{milcGris}{Saber contemporáneo}
La \textbf{iteración} es refinar un texto en varias versiones, cada una con un propósito distinto. En el oficio editorial se sabe hace mucho que el primer borrador nunca es el definitivo, y con IA generativa eso es más cierto, no menos. Una iteración típica tiene tres versiones. La \textbf{V1} es el borrador que sale del prompt profesional: captura estructura e ideas, pero suele traer lugares comunes y voz de chatbot. La \textbf{V2} es la revisión, donde le pides a la herramienta cambios concretos sobre ese texto. La \textbf{V3} es la versión con tu voz, y esa lleva escritura tuya directa. Sin eso, la V3 es una V2 disfrazada.
\end{softbox}

\begin{softbox}{milcMostaza}{milcGris}{Pregunta-puente}
\textit{¿Por qué un oficio entero se construye quitando hilos en vez de poniéndolos? ¿Y qué pasa cuando alguien entrega el primer borrador de una IA sin haberle quitado nada?}
\end{softbox}

\begin{softbox}{milcVerde}{milcGris}{Saber y saber hacer}
Al terminar podrás: (1) \textbf{identificar} en tu borrador qué material sirve y qué es relleno --- lugares comunes, repeticiones, frases vacías; (2) \textbf{aplicar} prompts de revisión concretos para producir una segunda versión mejor; (3) \textbf{evaluar} con honestidad cuánto de la versión final es tuyo y declararlo.
\end{softbox}



\small
\begin{tabularx}{\textwidth}{>{\bfseries}p{3.0cm}Y}
\rowcolor{milcGradePrimary!12}Autor & Dr. Álvaro Cárdenas Orozco\\
DBA & Producción editorial mediada por IA con criterio ético y técnico (MEN, T\&I)\\
Referentes & Ley 23 de 1982 · Floridi (ética de la IA generativa) · Dussel (autoría con dignidad)\\
Producto final & Un capítulo en 3 versiones progresivas, con bitácora de cambios, comparación lado a lado y una reflexión honesta sobre cuánto del texto final es tuyo\\
\end{tabularx}
\normalsize

\puente{En la sesión 3 conseguiste un primer borrador con un prompt profesional. Hoy lo conviertes en un capítulo. Entre esas dos cosas hay varias vueltas, y ninguna la hace sola la herramienta.}

\milcestacion{milcGradeDark}{Ruta MILC: así vamos a aprender}
\begin{tabularx}{\textwidth}{>{\centering\arraybackslash}X>{\centering\arraybackslash}X>{\centering\arraybackslash}X>{\centering\arraybackslash}X}
\phasecard{1}{milcVerde}{milcGris}{Escucha\\[-1mm]\small Observo, escucho y pregunto.} &
\phasecard{2}{milcTurquesa}{milcGris}{Entender\\[-1mm]\small Sistematizo y ordeno.} &
\phasecard{3}{milcMagenta}{milcGris}{Hacer\\[-1mm]\small Construyo y aplico.} &
\phasecard{4}{milcVino}{milcGris}{Cierre\\[-1mm]\small Evalúo y reflexiono.}
\end{tabularx}


\puente{Antes de teorizar, vas a leer tu propio borrador con lápiz en mano y a marcarlo. Solo cuando veas cuánto sobra vas a saber qué pedirle a la herramienta en la segunda vuelta.}

\milcestacion{milcVerde}{Estación 1 de 3 · Escucha}

\begin{softbox}{milcVerde}{milcGris}{Escena para escuchar}
\textbf{Actividad 1 · IDENTIFICA --- El borrador bajo la lupa} (15 min · individual). (1) Abre o imprime el borrador que generaste en la sesión 3. (2) Recórrelo entero marcando cada fragmento con una etiqueta: lugar común, repetición, relleno, dato sin verificar, voz ajena o material que sirve. (3) Cuenta cuántas marcas hay de cada tipo. (4) Calcula qué parte del texto se queda y qué parte hay que reescribir. (5) Anota la frase que más te molestó al releerla.
\end{softbox}

{\sffamily\bfseries\fontsize{8}{10}\selectfont\textcolor{milcNegro!55}{MARCA CUANDO LO HAYAS HECHO}}
\milccheckrow{Marqué el borrador completo, no solo los primeros párrafos.}
\milccheckrow{Conté cuántas marcas hay de cada tipo.}
\milccheckrow{Estimé qué parte del texto se queda y qué parte se reescribe.}

\begin{infoband}{milcVerde}


\textbf{Lo que va al cuaderno (Actividad 1 · IDENTIFICA).} \textbf{Título:} «Actividad 1 --- El borrador bajo la lupa». \textbf{Formato:} tabla con una fila por etiqueta y su conteo, la parte del texto que se queda y la frase que más te molestó. \textbf{Extensión:} un tercio de página. \textbf{Sabes que terminaste cuando} sabes qué pedirle a la herramienta en la segunda vuelta.
\end{infoband}

\puente{Ya marcaste el borrador. Ahora vas a entender la \textbf{anatomía de la iteración}: las tres versiones, cómo se escribe un prompt de revisión y qué errores hacen que la última versión no sea de nadie.}

\milcestacion{milcTurquesa}{Estación 2 de 3 · Entender}

\begin{softbox}{milcTurquesa}{milcGris}{Lo que vas a entender}
\textbf{Actividad 2 · APLICA --- Prompts de revisión} (20 min · parejas). (1) Con tu pareja, escriban las tres versiones con el propósito de cada una. (2) A partir de tus marcas, escribe entre tres y cinco prompts de revisión concretos. (3) Intercambien los prompts y marquen cuáles son vagos. (4) Reescriban juntos los vagos hasta que digan un cambio verificable.
\end{softbox}

\milcpaso{1}{milcTurquesa}{Las tres versiones se descomponen así. La \textbf{V1} sale de una sola pasada del prompt profesional: tu trabajo con ella es producirla y marcarla, no aprobarla. La \textbf{V2} nace de pedirle a la herramienta cambios concretos sobre ese texto. La \textbf{V3} lleva escritura tuya directa en los puntos donde la voz se nota.}
\milcpaso{2}{milcTurquesa}{Un prompt de revisión útil nombra un cambio verificable sobre un texto que le pegas. «Reescribe la primera oración como una pregunta» sirve. «Mejóralo» no sirve, porque no dice qué se va a comparar después. La diferencia entre los dos es si puedes decir, al leer la respuesta, si el cambio se hizo o no.}
\milcpaso{3}{milcTurquesa}{Todo se juega en una regla: una versión final sin mano humana directa no es final, es la segunda disfrazada. Los puntos donde más se nota la voz son la apertura, el cierre y los ejemplos propios. Si esos tres los escribió la herramienta, el capítulo va a sonar a herramienta por más vueltas que le des.}
\milcpaso{4}{milcTurquesa}{Algoritmo: (1) toma la V1 marcada; (2) escribe tres a cinco prompts de revisión; (3) genera la V2; (4) imprímela y toma lápiz; (5) reescribe a mano la apertura, el cierre y uno o dos ejemplos propios; (6) pasa a limpio la V3; (7) llena la bitácora con cada cambio y su razón.}

\begin{drawbox}{milcTurquesa}{Anatomía de la iteración}
\textbf{V1:} borrador de la herramienta con el prompt profesional. \\[1mm] \textbf{V2:} revisión con prompts concretos y verificables. \\[1mm] \textbf{V3:} voz propia, escrita a mano en los puntos críticos. \\[1mm] \textbf{Puntos críticos:} apertura, cierre y ejemplos propios. \\[1mm] \textbf{Regla:} sin mano humana no hay V3.
\end{drawbox}

\begin{softbox}{milcMostaza}{milcGris}{Errores comunes y su corrección}
(1) Entregar la V1 como definitiva, cuando es material crudo. (2) Saltarse el marcado y pasar directo a pedir mejoras. (3) Escribir prompts de revisión vagos, que no permiten comprobar nada. (4) Dejar pasar los lugares comunes: «en un mundo donde», «cada vez más», «en pleno siglo XXI». (5) Producir una V3 sin una sola línea escrita por ti.
\end{softbox}

\begin{infoband}{milcTurquesa}


\textbf{Lo que va al cuaderno (Actividad 2 · APLICA).} \textbf{Título:} «Actividad 2 --- Prompts de revisión». \textbf{Formato:} las tres versiones con su propósito, tus tres a cinco prompts de revisión y la corrección que hizo tu pareja de los que estaban vagos. \textbf{Extensión:} media página. \textbf{Sabes que terminaste cuando} cada prompt tuyo nombra un cambio que se puede comprobar.
\end{infoband}

\puente{Tienes el método. Toca aplicarlo. Generas la segunda versión con prompts de revisión, escribes a mano los puntos críticos de la tercera y documentas cada cambio.}

\milcestacion{milcMagenta}{Estación 3 de 3 · Hacer}

\begin{softbox}{milcMagenta}{milcGris}{Cómo vas a construir tu producto}
\textbf{Actividad 3 · EVALÚA --- Tres versiones y una reflexión} (40 min · individual). (1) Genera la V2 aplicando tus prompts de revisión. (2) Imprímela y reescribe a mano la apertura, el cierre y uno o dos ejemplos propios. (3) Pasa a limpio la V3 integrando lo que escribiste. (4) Llena la bitácora con versión, cambio, razón y resultado. (5) Arma la comparación de las tres versiones lado a lado. (6) Escribe cinco líneas sobre cuánto de este capítulo es tuyo, sin adornar la respuesta.
\end{softbox}

\milcpaso{1}{milcMagenta}{Tu trabajo se construye en tres piezas: las tres versiones del capítulo, la bitácora de cambios y la reflexión final. La bitácora es la que después te va a permitir iterar el capítulo 2 sin volver a esta guía.}
\milcpaso{2}{milcMagenta}{Sigue el patrón «mano humana en los puntos críticos». No hace falta reescribir todo a mano. Basta con que la apertura, el cierre y uno o dos ejemplos propios los hayas escrito tú. Son los lugares donde quien lee reconoce si hay alguien detrás.}
\milcpaso{3}{milcMagenta}{La bitácora tiene cuatro columnas: versión, cambio aplicado, razón del cambio y resultado. Al menos dos filas por versión. Lo que la hace útil es la columna de la razón, porque es la que se puede discutir.}
\milcpaso{4}{milcMagenta}{Algoritmo: (1) V2 con los prompts de revisión; (2) V3 con mano propia en los puntos críticos; (3) bitácora; (4) comparación lado a lado por fragmento; (5) reflexión de cinco líneas. Si al escribirla tienes que decir que casi todo es de la herramienta, eso también es un resultado y hay que anotarlo.}

\begin{drawbox}{milcMagenta}{Checklist antes de entregar el capítulo}
\checkbox Borrador V1 con las marcas de lectura crítica. \\[1mm] \checkbox V2 generada con prompts de revisión concretos. \\[1mm] \checkbox V3 con apertura, cierre y ejemplos escritos a mano. \\[1mm] \checkbox Bitácora con al menos dos filas por versión. \\[1mm] \checkbox Comparación de las tres versiones lado a lado. \\[1mm] \checkbox Reflexión de cinco líneas sin adornos.
\end{drawbox}

\begin{softbox}{milcMagenta}{milcGris}{Plantilla de trabajo}
\textbf{V1.} \textit{Borrador con marcas de lectura crítica.} \\[1mm] \textbf{V2.} \textit{Revisión con prompts concretos.} \\[1mm] \textbf{V3.} \textit{Voz propia en apertura, cierre y ejemplos.} \\[1mm] \textbf{Bitácora.} \textit{Versión, cambio, razón y resultado.} \\[1mm] \textbf{Comparación.} \textit{Las tres versiones, fragmento por fragmento.} \\[1mm] \textbf{Reflexión.} \textit{Cinco líneas sobre cuánto es tuyo.}
\end{softbox}

\begin{infoband}{milcMagenta}


\textbf{Lo que va al cuaderno (Actividad 3 · EVALÚA).} \textbf{Título:} «Actividad 3 --- Capítulo iterado». \textbf{Formato:} la referencia a las tres versiones, la bitácora, la comparación y la reflexión de cinco líneas. \textbf{Extensión:} media página. \textbf{Sabes que terminaste cuando} lees la V3 en voz alta y te suena a ti.
\end{infoband}

\puente{Lo que sigue resume \emph{qué} entregas hoy: las tres versiones, la bitácora de cambios y la reflexión. El criterio principal es que la versión final te suene a ti al leerla en voz alta.}

\milcestacion{milcMostaza}{Producto de la sesión}
\textbf{3 versiones del capítulo, bitácora de cambios, comparación y reflexión}.
\begin{softbox}{milcMostaza}{milcGris}{Criterios de calidad}
(1) Tres versiones progresivas del capítulo, con propósito distinto cada una. (2) Prompts de revisión que nombran cambios verificables. (3) Apertura, cierre y ejemplos escritos a mano en la versión final. (4) Bitácora con al menos dos cambios documentados por versión. (5) Comparación de las tres versiones fragmento por fragmento. (6) Reflexión honesta sobre cuánto del capítulo es propio.
\end{softbox}

\puente{Antes de cerrar, mira la iteración desde las cinco dimensiones humanas. Quitar lo que sobra de un texto propio cuesta, y esa dificultad es parte del oficio.}

\milcestacion{milcVino}{Evaluación liberadora · cinco dimensiones}

\begin{softbox}{milcVino}{milcGris}{Sentido de la evaluación}
Evaluar no es solo revisar una nota. Es reconocer qué aprendí, qué cambió en mí, qué emociones manejé, cómo me relaciono con la ciudad y la comunidad, y qué le dejaré a quien viene detrás.
\end{softbox}

\textbf{Mi reflexión liberadora en cinco dimensiones.} Completa cada frase iniciada:

\small
\begin{tabularx}{\textwidth}{>{\bfseries\RaggedRight\arraybackslash}p{3.4cm}Y>{\RaggedRight\arraybackslash}p{4.6cm}}
\rowcolor{milcVino}\color{white}Dimensión & \color{white}\bfseries Mi respuesta (frame starter) & \color{white}\bfseries Algo que descubrí\\
1. Personal & Lo que más cambió en mí fue \linead{6.5cm} & \linead{4.2cm}\\
2. Emocional & La emoción más fuerte fue \linead{2.5cm} y la regulé \linead{3cm} & \linead{4.2cm}\\
3. Ciudadana & En democracia esto importa porque \linead{6cm} & \linead{4.2cm}\\
4. Local (Cartago) & En mi barrio sería útil para \linead{6.5cm} & \linead{4.2cm}\\
5. Intergeneracional & A mi abuela/abuelo le diría \linead{6.5cm} & \linead{4.2cm}\\
\end{tabularx}
\normalsize

\begin{infoband}{milcVino}
\textbf{Para la próxima sustentación.} Vuelve a esta tabla en seis meses. Verás que las respuestas cambian — no porque lo aprendido cambie, sino porque tú cambias. La evaluación liberadora es una conversación contigo a través del tiempo.
\end{infoband}


\puente{Y para terminar, tres voces sobre volver sobre lo hecho. Dussel sobre el rodeo que hace posible la cercanía, Marco Aurelio sobre cambiar de parecer, Floridi sobre no saber cuál mitad sobra.}

\milcestacion{milcGradeDark}{Triángulo de pensamiento · cierre filosófico}

\begin{softbox}{milcVino}{milcGris}{Dussel · lente del nosotros}
\textit{«El maestro y el discípulo deben apartarse para preparar en la vida su discurso futuro… El rodeo de la lejanía hace posible la proximidad futura.»}

{\footnotesize\itshape\color{milcNegro!70}--- Enrique Dussel · Filosofía de la liberación (1977), §2.2.1.1}

Un texto recién escrito no se puede corregir bien. Hace falta apartarse: dejarlo, imprimirlo, leerlo en voz alta, verlo con otra letra. Ese rodeo es lo que permite leerlo como lo leería alguien más. Sin distancia, uno corrige comas y deja intacto lo que estaba mal desde el principio.

\textbf{¿Leí mi borrador como si fuera de otro, o lo leí defendiéndolo?}
\end{softbox}
\linead{\linewidth}\\[1mm]
\linead{\linewidth}

\begin{softbox}{milcMostaza}{milcGris}{Estoicismo · Marco Aurelio · Meditaciones VIII, 16 (c. 175 d.C.) · cuidado interior}
\textit{«Acuérdate que igualmente te es libre el mudar de parecer y el seguir el aviso de quien te corrija.»}

{\footnotesize\itshape\color{milcNegro!70}--- Marco Aurelio · Meditaciones VIII, 16 (c. 175 d.C.)}

Iterar es cambiar de parecer sobre lo que uno mismo escribió. La resistencia no viene de la herramienta: viene de que ya invertiste tiempo en esa versión. En el calado, quitar hilos es parte del trabajo y nadie lo vive como una pérdida.

\textbf{¿Qué fragmento defendí solo porque me costó escribirlo?}
\end{softbox}
\linead{\linewidth}\\[1mm]
\linead{\linewidth}

\begin{softbox}{milcTurquesa}{milcGris}{Floridi · lente de la infoesfera}
\textit{«Parafraseando un dicho de la publicidad: la mitad de nuestros datos es basura, solo que no sabemos cuál mitad. Lo que necesitamos es entender mejor qué datos vale la pena conservar. (trad. propia)»}

{\footnotesize\itshape\color{milcNegro!70}--- Luciano Floridi · Big data and their epistemological challenge (2012)}

Con un borrador generado pasa igual: buena parte sobra y el problema es saber cuál. Por eso la actividad 1 no es un trámite. Marcar el texto es lo que convierte una intuición vaga de que algo no está bien en una lista de cosas que se pueden pedir.

\textbf{¿Sabría señalar con el dedo qué mitad de mi borrador sobra, o solo siento que sobra algo?}
\end{softbox}
\linead{\linewidth}\\[1mm]
\linead{\linewidth}

\begin{infoband}{milcNegro}
\textbf{Nota para el docente.} Las tres son citas textuales y llevan su referencia debajo. Puedes remitir a la obra en clase.
\end{infoband}

\milcestacion{milcVino}{Mi autoevaluación · marca una casilla por fila}

{\small
\setlength{\extrarowheight}{2.2pt}
\begin{tabularx}{\textwidth}{>{\RaggedRight\arraybackslash\bfseries}p{3.0cm}YYY}
\rowcolor{milcVino}\color{white}\bfseries Criterio & \color{white}\bfseries Lo logré & \color{white}\bfseries En proceso & \color{white}\bfseries Necesito apoyo\\[.6mm]
Comprensión del tema & \milcopcion{Explico las ideas con mis palabras.} & \milcopcion{Reconozco las ideas, falta precisión.} & \milcopcion{Necesito repasar lo básico.}\\[.8mm]
\rowcolor{milcGris}Pensamiento computacional & \milcopcion{Usé los pilares al armar mi producto.} & \milcopcion{Usé algunos pilares.} & \milcopcion{Necesito releer la estación 2.}\\[.8mm]
Aplicación y producto & \milcopcion{Mi producto cumple los criterios.} & \milcopcion{Está armado, con ajustes pendientes.} & \milcopcion{Aún está en construcción.}\\[.8mm]
\rowcolor{milcGris}Reflexión 5 dimensiones & \milcopcion{Respondí cada una con honestidad.} & \milcopcion{Respondí algunas con detalle.} & \milcopcion{Mis respuestas son muy generales.}\\[.8mm]
Triángulo de pensamiento & \milcopcion{Conecté las 3 voces con mi trabajo.} & \milcopcion{Conecté al menos una.} & \milcopcion{No las relaciono con mi proceso.}\\[.8mm]
\end{tabularx}}

\vspace{3mm}
\textbf{Hoy entendí que:}\\[1mm]
\linead{\linewidth}\\[1mm]
\linead{\linewidth}

\textbf{Mi compromiso para la próxima sesión es:}\\[1mm]
\linead{\linewidth}\\[1mm]
\linead{\linewidth}

\begin{softbox}{milcMostaza}{milcGris}{Cita de despedida · Marco Aurelio}
\textit{``El obstáculo en el camino se vuelve el camino. Lo que estorba al camino se vuelve camino.''}\\[1mm]
Cada error de hoy, bien observado, se convierte en pieza del camino que pisarás mañana.
\end{softbox}

\small
\begin{tabularx}{\textwidth}{>{\bfseries}p{3.6cm}Y}
\rowcolor{milcGradeDark!12}Nombre completo: & \linead{10cm}\\
Fecha: & \linead{4cm} \quad Curso/grupo: \linead{4cm}\\
Mi firma: & \linead{9cm}\\
\end{tabularx}
\normalsize

\begin{infoband}{milcGradeDark}
\textbf{Para la próxima sesión.} Llega con la versión final del capítulo, la bitácora y la reflexión. En la sesión 5 vas a afinar la estructura completa del libro: el arco, el ritmo y el peso de cada capítulo. Lo que aprendiste iterando uno se aplica a los diez.
\end{infoband}

\end{document}
